\documentclass[10pt,a4paper]{article}

\usepackage[margin=1.05in]{geometry}
\usepackage{setspace}
\usepackage{amsmath,amssymb,amsthm}
\usepackage{graphicx}
\usepackage{booktabs}
\usepackage{array}
\usepackage{caption}
\usepackage{enumitem}
\usepackage{titlesec}
\usepackage{microtype}
\usepackage{hyperref}
\usepackage{fancyhdr}

\hypersetup{colorlinks=true,linkcolor=black,urlcolor=black,citecolor=black}
\captionsetup{font=small,labelfont=bf,skip=4pt,justification=raggedright}

\titleformat{\section}{\normalsize\bfseries}{}{0em}{}
\titleformat{\subsection}{\small\bfseries}{}{0em}{}
\titleformat{\subsubsection}{\small\itshape}{}{0em}{}
\titlespacing*{\section}{0pt}{1.35em}{0.5em}
\titlespacing*{\subsection}{0pt}{0.95em}{0.35em}

\theoremstyle{plain}
\newtheorem{proposition}{Proposition}
\newtheorem{lemma}{Lemma}
\theoremstyle{definition}
\newtheorem{definition}{Definition}
\theoremstyle{remark}
\newtheorem{observation}{Observation}

\setstretch{1.14}
\setlist[itemize]{leftmargin=1.15em,itemsep=0.1em,topsep=0.2em}

\graphicspath{{drawings/}}

\pagestyle{fancy}
\fancyhf{}
\fancyhead[L]{\small\itshape Elements of Position}
\fancyhead[R]{\small\itshape II.\ The Other Walk}
\fancyfoot[C]{\small\thepage}
\renewcommand{\headrulewidth}{0pt}

\begin{document}

\begin{center}
\vspace*{2.2em}
{\large Elements of Position}\\[1.6em]
{\LARGE\bfseries II.\ The Other Walk}\\[2.2em]
{\normalsize Justin Erholtz}\\[0.35em]
{\small 12 September 2026}
\end{center}

\vspace{2.4em}

\begin{center}
{\large 2 is named by 1.}
\end{center}

\vspace{1.8em}

\noindent
The second street is not a new unit --- it is the same unit asked the other way.

Part~I closed a hold and kept the cut in view, then repeated that hold with the $n$-count, integer by integer, at rate $1$.
This part does not reopen those proofs except to use them.
What it adds is a second rigid process on the same midline, and the observation that some arrivals do not factor through earlier finished counts.

$2$ is already named at count $1$, by inversion of the only primary cut, and walked at count $2$.
$\Phi=1/\varphi$ is written the same way $\pi=1/\pi$ was written: inverse parts of the unit, not matching decimals.
Both processes share the $n$-count, but only one walks out --- the other, after arrival, walks back toward the unit.

\vspace{1.4em}

\noindent
\textit{A note.}
This part cites \emph{I.\ The Closed Unit} and stays inside the same questions.
Claims that can be proved here are proved by school geometry.
The drawings are layers on the same architectural set, built by \texttt{Elements of Position - Drawings.py}.

\vspace{0.9em}
\noindent
\textit{A note on names used here.}
Names from I stand: closed unit, primary ratio, prime ratio, own-origin, street, unit, vacuum, trace, master sequence, $n$-count.
Other walk --- the second process, from a finished $N$ back toward the unit.
Similar cut --- $1/\varphi$, the other self-ratio on the same $1/2$ line; not a second primary ratio.
$\Phi$ --- the inverse of that cut, $\Phi=1/\varphi$.
Docking --- when the other walk lands on a count the $n$-count already holds.

\newpage
\tableofcontents
\newpage

% =====================================================================
\section{Named by one}
% =====================================================================

$1$ names $2$ --- not by walking to $2$ first, but by inverting the only primary cut of the finished unit $[0,1]$.

\begin{proposition}
\label{prop:two-named}
$2=1/(1/2)$.
On the pair $(x,2-x)$ the arithmetic mean is constantly $1$, and equals the geometric mean only at $x=1$.
At integer $n=2$,
\[
2\cdot\frac12=1,\qquad d=2,\qquad \frac{2\pi}{2}=\pi.
\]
\end{proposition}

\begin{proof}
The first sentence is inversion of the primary ratio already named in I.
The second is the fixed-sum case of AM--GM already proved there.
The third is the half-walk identity of I at $n=2$.
\end{proof}

So $1$ names $2$ additively as the next integer after the close, and $1$ holds $2$ by inversion.
The live mean of the diameter-$2$ walk is $\pi$ --- the first integer after the paper close, not a new constant.

At count $1$ the pair $\{1/2,\,2\}$ is named and the arithmetic mean of that pair is $5/4$.
Until count $2$ both ends of $[1/2,2]$ have not been walked, so $5/4$ is only a dated result.
At count $2$ the pair is built and $5/4$ becomes the ordinary midpoint of a walked segment.
The construction does not invent a location; it waits until both ends exist.

\subsection{The rate of the $n$-count}

The $n$-count is the master sequence $1,2,3,\ldots$, one unit per step, at a rate of unfold equal to $1$.
A run to a million, or ten million, does not change the time between integers.
What it changes is the depth of sampling: the local circle is cut into $n$ equal arcs, each still carrying the locked ratio $\pi$.

Rate $0$ is pause: one HERE, no next tick; any rate other than $0$ is the count moving.
A day between integers, a year, a millisecond --- no tempo reaches a different kind of mark.
The arrival package itself, $n$ and $1/n$ and the cross-ratio $-n$, does not care what watch you used.

% =====================================================================
\section{Irreducible arrivals}
% =====================================================================

An inversion through $1$ is a pair $(n,1/n)$ with product $1$, already proved in I for every positive $n$.
What this section notices is a grammar inside those pairs once the $n$-count has run far enough to have earlier letters on the table.

\begin{proposition}
\label{prop:reducible}
If $n=ab$ for integers $a,b>1$, then
\[
\frac1n=\frac1a\cdot\frac1b.
\]
The interior at $n$ is then a product of two earlier interiors.
\end{proposition}

\begin{proof}
Immediate from $n=ab$.
\end{proof}

\begin{proposition}
\label{prop:irreducible}
If no such $a,b$ exist, $1/n$ cannot be written from any earlier interior of the count.
Those $n$ are the primes already arrived at.
The first such pair is $(2,1/2)$.
\end{proposition}

\begin{proof}
That is the definition of a prime among the integers already on the table.
$2$ is the first integer after $1$ that is not $1$, and $2$ has no factorization $a\cdot b$ with $a,b>1$.
\end{proof}

The observation comes first.
The label ``prime'' comes after.
There is no completed prime sequence in this construction, because the $n$-count is not complete.
What there are, are the irreducible inversions already arrived at; before arrival there is nothing there to name.

$1/2$ is doing two jobs at once: it is the even split of any segment --- the AM--GM cut of the unit, the primary ratio --- and it is also the interior of the first irreducible inversion.
Every later diameter still carries that cut as the midpoint $n/2$.

Composites are words in earlier letters: $1/4=(1/2)^{2}$, $1/6=(1/2)(1/3)$.
At a high end-count the grammar does not change; the clock does not wait for a finished list.

The naming choice from I stays in force: that $2$ is both the inverse of the primary cut and the first prime is a coincidence of small numbers worth marking, not a claim that primality follows from inversion.
No later result in this part depends on that convergence holding for any $n$ other than $2$.
What generalizes is the observation of irreducibility after arrival.

\begin{figure}[ht]
\centering
\includegraphics[width=0.95\textwidth]{layer_primes_rate.png}
\caption{The rate of the $n$-count remains $1$.
Some $n$ refuse to factor from earlier interiors.}
\end{figure}

\subsection{Unitization last}

Classically a measurement is a number and then a unit: fourteen, then kilograms.
Here the order is the reverse of that habit, and it is forced.

Arrive, then measure as a ratio of parts already on the table, then --- if at all --- call the closed $1$ a unit.
No imported stick is used.
If something is a unit, it is $1$ measuring itself.

Nothing discussed is an imported completed object, and no sequence found within is complete: the primes on the table are the refusals so far, the Fibonacci integers on the table are the arrivals so far.
$\Phi$ is a cut of the unit, not a warehouse of later copies.

% =====================================================================
\section{The other similar cut}
% =====================================================================

On the same midline there is another cut that names itself by inversion.
It is not the primary ratio, but a similar cut of a different kind.

\begin{proposition}
\label{prop:phi}
The equation
\[
\frac1x=1+x
\]
has a unique root in $(0,1)$,
\[
x=\frac{\sqrt5-1}{2}=\Phi:=\frac1\varphi,\qquad
\varphi=\frac{1+\sqrt5}{2}.
\]
Equivalently $\varphi=1+\Phi$ and $\varphi\cdot\Phi=1$.
\end{proposition}

\begin{proof}
$1/x=1+x$ is $x^{2}+x-1=0$.
The positive roots are $\varphi$ and $\Phi=\varphi-1=1/\varphi$.
Only $\Phi$ lies in $(0,1)$.
\end{proof}

That equation is the similar-cut condition on a unit segment: the part that remains is to the cut as the cut is to the whole.
The word \emph{ratio} is doing the same work it already did in I: $\varphi$ is whole over larger part, $\Phi$ is larger part over whole.
No other name is required, and the word golden is not used.

\begin{proposition}
\label{prop:phi-gm}
\[
\varphi\cdot\Phi=1,\qquad
\mathrm{GM}(\varphi,\Phi)=1,\qquad
\mathrm{AM}(\varphi,\Phi)=\frac{\sqrt5}{2}>1.
\]
\end{proposition}

\begin{proof}
The product is $1$ by Proposition~\ref{prop:phi}.
The geometric mean is $\sqrt{1}=1$.
The arithmetic mean is $(\varphi+\Phi)/2=\sqrt5/2$, which exceeds $1$ because $\varphi\neq\Phi$.
\end{proof}

So $\varphi$ and $\Phi$ are an inversion pair through $1$, of the same kind as $(n,1/n)$ and $(\pi,1/\pi)$.
They are not an integer pair --- they are the unique pair whose larger part is the smaller plus the unit.

Equal, here, still does not mean the decimals match.
It means the parts name each other through the whole.

\subsection{They do not leave the even cut}

\begin{proposition}
\label{prop:midline}
On the unit segment,
\[
\Phi+\Phi^{2}=1,\qquad
\Phi-\frac12=\frac12-\Phi^{2}.
\]
The two parts sit equally far from the midpoint $1/2$.
\end{proposition}

\begin{proof}
From $\varphi^{2}=\varphi+1$ divided by $\varphi^{2}$ one has $1=\Phi+\Phi^{2}$, so $\Phi^{2}=1-\Phi$.
The two distances to $1/2$ are then $\Phi-1/2$ and $1/2-\Phi^{2}=\Phi-1/2$.
\end{proof}

That is the same hinge as I, read on the segment.
$\pi$ and $1/\pi$ do not leave the unit; their geometric mean is $1$.
$\varphi$ and $\Phi$ do not leave the unit; their geometric mean is $1$.
On the paper diameter they additionally do not leave the even cut: $\Phi$ and $\Phi^{2}$ revolve around $1/2$.

Where is $\varphi$?
$\Phi$ is the interior mark; $\varphi$ is the outward similar copy of the paper circle --- diameter $\varphi$, same centre, same $1/2$ midline.
The paper circle is the unit itself.
Inward, one similar copy has diameter $\Phi$; outward, one has diameter $\varphi$.
They do not leave each other.

\begin{figure}[ht]
\centering
\includegraphics[width=0.95\textwidth]{layer_phi_10_midline_and_pair.png}
\caption{$\Phi$ and $\Phi^{2}$ on the paper diameter, symmetric about $1/2$.
On $xy=1$, the pair $(\varphi,\Phi)$ sits with $(1,1)$, $(2,1/2)$, and $(\pi,1/\pi)$.}
\end{figure}

\begin{figure}[ht]
\centering
\includegraphics[width=0.52\textwidth]{layer_phi_11_where_phi_is.png}
\caption{Where $\varphi$ is: concentric diameters $\Phi$, $1$, $\varphi$.}
\end{figure}

% =====================================================================
\section{Two processes, one $n$-count}
% =====================================================================

The generator carries two lines: they share the $n$-count, but they do not share a direction.

\paragraph{Process A.}
The body of I.
Paper turn at $n=1$, then radius $n/2$, one turn per integer, half-walk $n\pi/2$, inscribed $n$-gon.
Walked out with the count.

\paragraph{Process B.}
The other walk.
Same integers.
After arrival at a finished $N$, opposite angle, radius scaled by $\Phi$ at each opposite turn, back toward the unit.

\begin{proposition}
\label{prop:two-proc}
Let $N$ be a finished integer.
Process A at parameter $n\in[1,N]$ has radius $n/2$.
Process B at depth $d\in[0,N-1]$ after that arrival has
\[
r=\frac N2\,\Phi^{d},\qquad
\theta=-2\pi d.
\]
At $d=0$ both processes sit on the circle of diameter $N$.
\end{proposition}

\begin{proof}
A is the walk of I, conical in the count.
B is the similar copy of that arrival circle, scaled by $\Phi$ at each opposite turn.
At $d=0$, $r=N/2$, which is the local circle of A at $n=N$.
\end{proof}

You can construct the similar copies as $n$ grows, but you cannot walk B until you have arrived, because the walk is the other way.
Vacuum may name the nest in one stroke; trace must wait for the integer.

Fibonacci integers already on the table --- $1,1,2,3,5,8,13$ inside a present count --- are arrivals at which consecutive integers already built can be measured against each other.
Their ratios lean toward $\varphi$ --- an observation after arrival, not a completed sequence brought in from outside.

\begin{figure}[ht]
\centering
\includegraphics[width=0.95\textwidth]{layer_both_01_paper.png}
\caption{Both processes on the paper circle.
The $1/2$ midline; $\Phi$ and $\Phi^{2}$; half-walk; inward copies.}
\end{figure}

\begin{figure}[ht]
\centering
\includegraphics[width=0.95\textwidth]{layer_both_02_interiors.png}
\caption{Local interiors.
A: $n$-gon and half-walk.
B: first inward copy at scale $\Phi$, and the pair $\Phi,\Phi^{2}$ on the diameter.}
\end{figure}

\begin{figure}[ht]
\centering
\includegraphics[width=0.78\textwidth]{layer_both_03_plan.png}
\caption{Plan.
A walks out with the $n$-count.
B is the other walk from arrival, opposite angle, scale $\Phi$.}
\end{figure}

\begin{figure}[ht]
\centering
\includegraphics[width=0.95\textwidth]{layer_both_04_obliques.png}
\caption{Obliques of the two lines.
A up with the count.
B down from arrival.}
\end{figure}

\begin{figure}[ht]
\centering
\includegraphics[width=0.72\textwidth]{layer_both_05_oblique.png}
\caption{Single oblique.
The two processes share the paper close and then part.}
\end{figure}

The question was never just where $1$ sits --- it was whether a finished thing can still be cut.
Process B is that inward face made into a walk, after arrival, without declaring the nest complete.

% =====================================================================
\section{What I already said, read through this layer}
% =====================================================================

Part~I began at a tension: $1$ as a close, and a cut between $0$ and $1$ that does not end.

The body held $\pi$ and $1/\pi$ as inverse parts of the paper close.
This part holds $\varphi$ and $\Phi$ the same way, and shows they sit on the even cut $1/2$ as well.

The body refused a completed limit as $n$ grows.
This part refuses a completed list of primes and a completed Fibonacci sequence for the same reason: observation after arrival.

The body treated scale as ratio; Process B is that sentence pointed the other way --- each similar copy is $\Phi$ of the last, and $\varphi$ of the next.

The cone of slope $1/2$ in I is the primary ratio in elevation; the midline of $\Phi$ and $\Phi^{2}$ is the primary ratio on the diameter.
Same $1/2$.

Meetings of A's spokes with B's marks are already visible at small $n$.
They are not proved here.
They are III.

% =====================================================================
\section{What stands}
% =====================================================================

Forced, and available once the $n$-count has begun:
$2=1/(1/2)$, named at count $1$ and walked at count $2$;
rate of the $n$-count equal to $1$, independent of tempo;
reducible interiors as products of earlier interiors;
irreducible arrivals as those $n$ that do not factor through earlier finished counts;
$\Phi=1/\varphi$ the unique root of $1/x=1+x$ in $(0,1)$;
$\varphi\cdot\Phi=1$, $\mathrm{GM}(\varphi,\Phi)=1$;
$\Phi$ and $\Phi^{2}$ equally far from $1/2$;
two processes on one $n$-count, B walked only after arrival.

Named, so they can be pointed at:
other walk, similar cut, $\Phi$ as inverse parts, docking.

Not granted:
a completed prime list;
$\varphi$ as a constant carried in from outside the cut that produced it;
an other walk that exists as a trace before the integer it starts from;
a second primary ratio --- there is still only $1/2$.

The first sentence of I is still the sentence.
Where is $1$?
$2$ is one answer already named.
The other walk is the same $1$, asked back toward itself.

\end{document}
